\section{Canonical heat evolution and the real initial-data completion}
\label{sec:source}
\subsection{Normalization and exact evolution}
Use the classical normalization \cite{polymath}
\begin{subequations}\label{eq:canonical}
\begin{align}
 \Phi(v)&=\sum_{n=1}^{\infty}
 (2\pi^2n^4\e^{9v}-3\pi n^2\e^{5v})\e^{-\pi n^2\e^{4v}},\\
 H(t,z)&=\int_0^\infty \e^{tv^2}\Phi(v)\cos(zv)\dd v,
 & H(0,z)&=\frac18\xi\left(\frac{1+\ii z}{2}\right).
\end{align}
\end{subequations}
The source is integrable and decays faster than a Gaussian.
Differentiation gives \(H_t=-H_{zz}\).
We write \(h(T,z)=H(-T,z)\), so \(h_T=h_{zz}\).
Let
\[
 k_T(v)=\frac1{\sqrt{4\pi T}}\exp\left(-\frac{v^2}{4T}\right),
 \qquad T>0.
\]
\begin{lemma}[Real initial-value formula]\label{lem:heat}
For every \(T>0\) and \(z\in\C\),
\begin{equation}
 h(T,z)=\int_\R k_T(z-u)H(0,u)\dd u.
 \label{eq:heat-convolution}
\end{equation}
The integral and its spatial and time derivatives converge uniformly
on compact subsets with \(T>0\).
\end{lemma}
\begin{proof}
Insert the cosine transform of \(H(0,u)\).
For real \(u\), the cosine is bounded, while
\[
 \abs{k_T(z-u)}
 =\frac{\e^{(\Ima z)^2/(4T)}}{\sqrt{4\pi T}}
      \e^{-(u-\Rea z)^2/(4T)}.
\]
Integrability of \(\Phi\) permits Fubini. The Fourier transform of
the Gaussian is \(\e^{-Tv^2+\ii zv}\), giving
\eqref{eq:heat-convolution}. Derivatives introduce fixed polynomial
factors and inverse powers of \(T\), still integrable on the stated
compact sets.
\end{proof}
This identity is the reason real Riemann--Siegel data can be used
without asserting a complex remainder theorem: the heat kernel
produces the holomorphic field.

\subsection{Charts, gamma scale, and frozen data}
Fix \(r,b>0\), \(0<\tau\le B<\infty\), and
\[
 \Omega_X(r,b)=\{z:\abs{\Rea z-X}\le r,\ \abs{\Ima z}\le b\},
 \quad R=r+1,\quad Y=b+1,\quad W=R+Y.
\]
The enlarged chart is \(\Omega_X(R,Y)\).
For physical applications the inherited choice is
\(r=\ell+2\), \(b=1+h_{\mathrm{av}}+2\); no physical estimate is
deduced here.
Throughout put
\[
 L=\log(X/(4\pi)),\quad a_X=\sqrt{X/(4\pi)},\quad
 N=\lfloor a_X\rfloor,\quad
 \bar B=\max(1,B),\quad V_B=8\sqrt{\bar B}\,L.
\]
The center \(X\), the orders, all coefficients, all real partitions,
and the endpoints \(X\pm V_B\) are fixed under both \(z\)- and
\(T\)-differentiation.

Define
\begin{subequations}\label{eq:gamma-scale}
\begin{align}
 w_\epsilon(z)&=(1+\epsilon\ii z)/2,\qquad \epsilon=\pm1,\\
 M(w)&=\frac{w(w-1)}{16}\pi^{-w/2}\Gamma(w/2),\\
 G(z)&=\{M(w_+(z))M(w_-(z))\}^{1/2},\qquad g=G(X)>0.
\end{align}
\end{subequations}
The square root is positive on the real segment and holomorphic
on sufficiently high fixed charts. Sectorial gamma estimates
\cite{dlmfgamma} give
\begin{equation}
 \frac{G'(z)}{G(z)}=-\frac{\pi}{8}+O(X^{-1}),\qquad
 G(X)\ge g_-X^{7/4}\e^{-\pi X/8},\qquad
 C_\chi\ge\sup_{\Omega_X(r,b)}\frac{g}{\abs{G(z)}},
 \label{eq:gamma-comparison}
\end{equation}
with fixed positive constants after a fixed onset.
For the growing-order slab, \cref{sec:calibration} supplies their explicit
values, full complex domains, and a common sufficient center onset.
Here and below classical constants may depend on fixed sectors
and chart radii, but not on an order subsequently allowed to grow.

For a field \(E\), the natural comparison norm is
\begin{equation}
 \normg{E}=\sup_{\substack{\tau\le T\le B\\z\in\Omega_X(r,b)}}
 \frac{\max_{0\le j\le3}\abs{\partial_z^jE(T,z)}
                      +\abs{\partial_TE(T,z)}}{\abs{G(z)}}.
 \label{eq:natural-norm}
\end{equation}
The center norm replaces \(\abs{G(z)}\) by \(g\); later the outer
norm replaces it by a positive amplitude envelope.
None divides by the field itself.

\subsection{Classical real Riemann--Siegel input}
\label{subsec:rs-input}
For real \(u\) in \(I_{X,B}=[X-V_B,X+V_B]\), write
\[
 a(u)=\sqrt{u/(4\pi)},\quad m(u)=\lfloor a(u)\rfloor,\quad
 \chi(u)=1-2(a(u)-m(u)),
\]
\[
 U(u)=\exp\left[-\ii\left\{\frac u4\log\frac u{4\pi}
                                   -\frac u4-\frac{\pi}{8}\right\}\right].
\]
The fractional coordinate \(\chi\) is distinct from the fluctuation
order \(p\). Set
\begin{subequations}\label{eq:real-completion}
\begin{align}
 S_m(u)&=\sum_{\epsilon=\pm1}\sum_{n=1}^m
                 M(w_\epsilon(u))n^{-w_\epsilon(u)},\\
 R_{m,K}(u)&=2\Rea\left\{(-1)^{m-1}M(w_+(u))U(u)a(u)^{-1/2}
                 \sum_{k=0}^{K}C_k(\chi(u),1/2)a(u)^{-k}\right\},\\
 f_K(u)&=S_{m(u)}(u)+R_{m(u),K}(u).
\end{align}
\end{subequations}
The coefficient normalization is the complex one in Polymath,
Proposition 6.2, equations (53)--(56) \cite{polymath}.
At \(\sigma=1/2\), for every positive integer \(K\), its
real initial-data remainder gives
\begin{equation}
 \abs{H(0,u)-f_K(u)}
 \le\frac{2^{7/4}}7G(u)a(u)^{-1/2}
       \Gamma((K+1)/2)\left(\frac{1.1}{a(u)}\right)^{K+1}.
 \label{eq:rs-remainder}
\end{equation}
Only this real initial formula is imported. No positive-time
evolved remainder is continued to negative time.

The complete change of cutoff is retained by setting
\begin{subequations}\label{eq:retained-integral}
\begin{align}
 Q_{X,K}(u)&=S_{m(u)}(u)-S_N(u)+R_{m(u),K}(u),\\
 \mathcal I_{X,K,B}(T,z)&=\int_{I_{X,B}}k_T(z-u)Q_{X,K}(u)\dd u.
\end{align}
\end{subequations}
If \(a(X+V_B)-a(X-V_B)<1\), there is at most one threshold
\(4\pi n^2\) in the window, and at most two real branches.
Their endpoints are independent of \(z,T\). Possible values at a
threshold have measure zero. The real-part operation remains inside
the initial data in \eqref{eq:real-completion}; moving it outside
the complex-kernel integral would change the function.

For completeness, the classical real coefficient estimate
\(\abs{C_k}\le (3\sqrt2/(2\pi))\Gamma(k/2)/2^k\), \(k\ge1\),
and \(\abs{C_0}\le1/2\), imply a bounded retained coefficient sum
when \(K\le a_{\min}^2\), where
\(a_{\min}^2=(X-V_B)/(4\pi)\). Indeed successive terms of each
parity have ratio \(k/(8a(u)^2)\le1/8\).
Together with the at-most-one exact transfer this gives
\(\abs{Q_{X,K}(u)}\le C G(u)X^{-1/4}\).
This is a size bound, not a regularity or quadrature theorem.

\subsection{The smaller outer envelope}
Let \(C_\pm(t,z)=A_{\mathrm o}(t,z)\e^{\pm\ii\Theta_{\mathrm o}(t,z)}/2\)
and \(\calE_{\mathrm o}(t,z)=\abs{C_+}+\abs{C_-}\).
For \(1\le T\le2\), gamma calibration yields
\begin{equation}
 \calE_{\mathrm o}(-T,z)\ge
 c_{\mathrm o}\abs{G(z)}
    \exp\left(-\frac{TL^2}{16}+\frac{\pi^2T}{64}\right).
 \label{eq:outer-calibration}
\end{equation}
One way to see the scale is to compare the modulus of
\(A_{\mathrm o}(0,z)\) with \(\abs{G(z)}\) by sectorial Stirling.
On a fixed chart \(L(z)=L+O(X^{-1})\), so the heat factor changes
by \(\exp(O((1+L)/X))\). Finally
\((\abs{\e^{\ii\Theta}}+\abs{\e^{-\ii\Theta}})/2\ge1\).
After increasing the onset, fixed comparison constants give
\eqref{eq:outer-calibration}, as in \cite{ans03,ans05}.
Thus
\begin{equation}
 \normo{E}\le c_{\mathrm o}^{-1}\e^{L^2/8}
       \lVert E\rVert_{\gamma;X;1,2}^{(3,1)}.
 \label{eq:outer-conversion}
\end{equation}
The factor \(\e^{L^2/8}\) is why fixed algebraic error estimates
alone do not establish the desired outer accuracy.
